\documentclass[10pt,a4paper]{amsart}
\usepackage[margin=19mm]{geometry}
\usepackage{amsmath,amssymb,mathtools}
\usepackage[scr=rsfs,cal=euler]{mathalfa}
\usepackage{xfrac}
\usepackage[unicode,hidelinks,bookmarks=false]{hyperref}
\newcommand{\ie}{i.e.\ }
\newcommand{\cf}{cf.\ }
\newcommand{\resp}{respectively\ }
\newcommand{\Subsection}{\subsection}
\providecommand{\nobreakdash}{\nobreak\hbox{-}}
\setlength{\emergencystretch}{2em}
\multlinegap0pt
\usepackage{amssymb}
\usepackage{mathtools}
\usepackage{float}
\usepackage{tikz}
\newtheorem{lemma}{Lemma}
\title{Equivalence of Almgren--Pitts and phase-transition half-volume spectra}
\author{Talant Talipov}
\date{Source: arXiv:2604.01091v1 (CC BY 4.0)}
\begin{document}
\maketitle

\noindent\emph{Source and licence.} Equivalence of Almgren--Pitts and phase-transition half-volume spectra. Version arXiv:2604.01091v1 (CC BY 4.0), distributed by arXiv under the Creative Commons Attribution 4.0 International licence, http://creativecommons.org/licenses/by/4.0/, as recorded in the arXiv licence metadata for this exact version and shown on the abstract page https://arxiv.org/abs/2604.01091v1. Full licence notice: \texttt{licenses/arxiv-2604.01091-CC-BY-4.0.txt}.\par\smallskip

\noindent\textit{Editorial scope.} Counted unit: the article's lemma lem:meanzero\_near\_half\_volume (source0.tex lines 605--611). The article's complete original proof of it is delivered with it (source0.tex lines 613--741, one \texttt{proof} environment of 129 lines). No mathematics is written, repaired or re-derived: the statement and the proof are the article's own lines, in the article's own notation, and no retained line is substituted or re-flowed.  No further source-local context is needed: every cross-reference of every retained line resolves inside the two delivered spans.  Nothing was omitted from any retained span, so the package carries no omission record.  No retained line contains a citation, so the wrapper carries no bibliography and no bibliography entry.  No retained line contains a figure environment, an image-inclusion command or drawing code, so no image is delivered and the package declares no asset dependency.  Editorial formatting only, in addition: an AMS \texttt{amsart} preamble that restates the article's own declarations and macros used by the retained lines, namely the environments \texttt{lemma} on separate counters, together with the standard packages those lines need.  The retained environments print in this document as Lemma 1 in order of appearance, whereas the article prints the same items with the numbering of its own full document.  The complete native source of the article as distributed in the arXiv e-print for this exact version is bundled unchanged as \texttt{sources/197637/source0.tex}.
\begin{lemma}\label{lem:meanzero_near_half_volume}
Fix $\delta\in(0,1)$ and let $q>1$, $C_1>0$, $\beta>1$ be as in {\rm (vii)} for the potential $W$. There exist constants $C_0=C_0(\delta)>0$ and $C_1'=C_1'(\delta)>0$ such that the following holds. If $u\in H^{1}(M) \setminus \{0\}$ satisfies $\int_M u =0$, then for every $t\in(-1+\delta,\,1-\delta)$,
\[
\Bigl|\operatorname{Vol}\bigl(\{u>t\}\bigr)-\mathfrak h\Bigr|
\ \le\ C_0\,\varepsilon\,E_\varepsilon(u)\ +\ C_1'\,\bigl(\varepsilon E_\varepsilon(u)\bigr)^{1/q}.
\]
\end{lemma}

\begin{proof}
Let
\[
\begin{aligned}
B_\delta(u)
&:=\Bigl\{x\in M:\ 
\min\bigl\{|u(x)-1|,\ |u(x)+1|\bigr\}\ge \delta\Bigr\},\\
A^\pm_\delta(u)
&:=\Bigl\{x\in M:\ |u(x)\mp 1|<\delta\Bigr\}.
\end{aligned}
\]
Then $M=A^+_\delta(u)\sqcup A^-_\delta(u)\sqcup B_\delta(u)$.

\smallskip
\noindent\emph{Claim 1.} There exists $m_\delta=m_\delta(\delta)>0$ such that
$W(s)\ge m_\delta$ whenever $\min\{|s-1|,|s+1|\}\ge\delta$. In particular,
\begin{equation}\label{eq:vol_Bdelta_bound}
\operatorname{Vol}\bigl(B_\delta(u)\bigr)\le \frac{\varepsilon}{m_\delta}\,E_\varepsilon(u).
\end{equation}
Indeed, using {\rm (vii)} we have $W(s)\ge C_1\beta^q$ for $|s|\ge \beta$, while on the compact set
$\{\,|s|\le\beta:\ \min\{|s-1|,|s+1|\}\ge\delta\,\}$ the continuous function $W$ attains a positive minimum
because $W>0$ away from $\{\pm1\}$. Taking $m_\delta$ to be the minimum of these two lower bounds gives the claim,
and \eqref{eq:vol_Bdelta_bound} follows from
\[
E_\varepsilon(u)\ \ge\ \int_{B_\delta(u)}\frac{W(u)}{\varepsilon}\,dV\ \ge\ \frac{m_\delta}{\varepsilon}\,\operatorname{Vol}\bigl(B_\delta(u)\bigr).
\]

\smallskip
Fix $t\in(-1+\delta,1-\delta)$ and set $\Omega_t:=\{u>t\}$.

\smallskip
\noindent\emph{Claim 2.} One has
\begin{equation}\label{eq:Omega_vs_Aplus_clean}
A^+_\delta(u)\subset \Omega_t\subset A^+_\delta(u)\cup B_\delta(u),
\qquad
\Omega_t\cap A^-_\delta(u)=\emptyset,
\end{equation}
and therefore
\begin{equation}\label{eq:Omega_vs_Aplus_vol}
\bigl|\operatorname{Vol}(\Omega_t)-\operatorname{Vol}(A^+_\delta(u))\bigr|\le \operatorname{Vol}\bigl(B_\delta(u)\bigr).
\end{equation}
Indeed, $u\ge 1-\delta>t$ on $A^+_\delta(u)$, so $A^+_\delta(u)\subset \Omega_t$.
Also $u\le -1+\delta<t$ on $A^-_\delta(u)$, so $\Omega_t\cap A^-_\delta(u)=\emptyset$.
Since $M=A^+_\delta(u)\sqcup A^-_\delta(u)\sqcup B_\delta(u)$, this implies
$\Omega_t\subset A^+_\delta(u)\cup B_\delta(u)$, proving \eqref{eq:Omega_vs_Aplus_clean}.
Taking volumes gives \eqref{eq:Omega_vs_Aplus_vol}.

\smallskip
\noindent\emph{Claim 3.} One has
\begin{equation}\label{eq:Aplus_half_clean}
\Bigl|\operatorname{Vol}\bigl(A^+_\delta(u)\bigr)-\mathfrak h\Bigr|
\le \frac12\,\operatorname{Vol}\bigl(B_\delta(u)\bigr)+\frac{1}{2(1-\delta)}\int_{B_\delta(u)}|u|\,dV.
\end{equation}
Indeed, since $\int_M u=0$,
\[
0=\int_{A^+_\delta(u)}u+\int_{A^-_\delta(u)}u+\int_{B_\delta(u)}u.
\]
Using $u\ge 1-\delta$ on $A^+_\delta(u)$ and $u\le -(1-\delta)$ on $A^-_\delta(u)$, we obtain
\begin{equation}
(1-\delta)\,\bigl(\operatorname{Vol}(A^+_\delta(u))-\operatorname{Vol}(A^-_\delta(u))\bigr)
\le \int_{B_\delta(u)}|u|\,dV.
\end{equation}
Note that $\int_M (-u)=0$ and that
$A^+_\delta(-u)=A^-_\delta(u)$, $A^-_\delta(-u)=A^+_\delta(u)$. Applying the same argument to $-u$ yields the reverse inequality
\[
(1-\delta)\,\bigl(\operatorname{Vol}(A^-_\delta(u))-\operatorname{Vol}(A^+_\delta(u))\bigr)
\le \int_{B_\delta(u)}|u|\,dV.
\]
Hence,
\begin{equation}\label{eq:Aplus_minus_Aminus_abs}
\bigl|\operatorname{Vol}(A^+_\delta(u))-\operatorname{Vol}(A^-_\delta(u))\bigr|
\le \frac{1}{1-\delta}\int_{B_\delta(u)}|u|\,dV.
\end{equation}
Since $\operatorname{Vol}(A^+_\delta(u))+\operatorname{Vol}(A^-_\delta(u))=\operatorname{Vol}(M)-\operatorname{Vol}(B_\delta(u))$, we can write
\[
\operatorname{Vol}(A^+_\delta(u))-\mathfrak h
=\frac12\bigl(\operatorname{Vol}(A^+_\delta(u))-\operatorname{Vol}(A^-_\delta(u))\bigr)-\frac12\operatorname{Vol}(B_\delta(u)),
\]
and taking absolute values and using \eqref{eq:Aplus_minus_Aminus_abs} yields \eqref{eq:Aplus_half_clean}.

\smallskip
We now estimate $\int_{B_\delta(u)}|u|$. Splitting into bounded and unbounded parts,
\[
\int_{B_\delta(u)}|u|
\le \int_{B_\delta(u)\cap\{|u|\le \beta\}}|u|+\int_{\{|u|>\beta\}}|u|
\le \beta\,\operatorname{Vol}\bigl(B_\delta(u)\bigr)+\int_{\{|u|>\beta\}}|u|.
\]
By H\"older and {\rm (vii)},
\begin{align*}
\int_{\{|u|>\beta\}} |u|
&\le \operatorname{Vol}(M)^{1-1/q}\Bigl(\int_{\{|u|>\beta\}} |u|^q\Bigr)^{1/q}\\
&\le \operatorname{Vol}(M)^{1-1/q}\Bigl(\frac{1}{C_1}\int_M W(u)\Bigr)^{1/q}\\
&\le \operatorname{Vol}(M)^{1-1/q}\Bigl(\frac{\varepsilon}{C_1}E_\varepsilon(u)\Bigr)^{1/q}.
\end{align*}
Insert this bound into \eqref{eq:Aplus_half_clean} to obtain
\[
\begin{aligned}
\Bigl|\operatorname{Vol}\bigl(A^+_\delta(u)\bigr)-\mathfrak h\Bigr|
&\le \frac12\,\operatorname{Vol}\bigl(B_\delta(u)\bigr)
+\frac{1}{2(1-\delta)}
\Biggl[
\beta\,\operatorname{Vol}\bigl(B_\delta(u)\bigr)\\
&\qquad\qquad
+ \operatorname{Vol}(M)^{1-1/q}\Bigl(\frac{\varepsilon}{C_1}E_\varepsilon(u)\Bigr)^{1/q}
\Biggr].
\end{aligned}
\]
Finally combine this with \eqref{eq:Omega_vs_Aplus_vol} and use the triangle inequality:
\[
\Bigl|\operatorname{Vol}(\Omega_t)-\mathfrak h\Bigr|
\le \bigl|\operatorname{Vol}(\Omega_t)-\operatorname{Vol}(A^+_\delta(u))\bigr|
     +\Bigl|\operatorname{Vol}(A^+_\delta(u))-\mathfrak h\Bigr|
\]
\[
\le \operatorname{Vol}(B_\delta(u))
+\frac12\,\operatorname{Vol}(B_\delta(u))
+\frac{\beta}{2(1-\delta)}\,\operatorname{Vol}(B_\delta(u))
+\frac{\operatorname{Vol}(M)^{1-1/q}}{2(1-\delta)}\Bigl(\frac{\varepsilon}{C_1}E_\varepsilon(u)\Bigr)^{1/q}.
\]

\smallskip
To conclude, use \eqref{eq:vol_Bdelta_bound} to absorb $\operatorname{Vol}(B_\delta(u))$ into $\varepsilon E_\varepsilon(u)$, and set
\[
C_0:=\left(\frac32+\frac{\beta}{2(1-\delta)}\right)\frac{1}{m_\delta},
\qquad
C_1':=\frac{\operatorname{Vol}(M)^{1-1/q}}{2(1-\delta)}\,C_1^{-1/q}.
\]
This gives the stated estimate.
\end{proof}

\end{document}
